\section{Method}
We propose \name, a decoupled \yl{Gaussian} splatting representation that decomposes geometry, texture, and lighting from multi-view images, which enables downstream applications like scene relighting and texture editing. 
The overview of our method is shown in Fig.~\ref{fig:pipeline}. 
First, we briefly introduce the rendering formulation in \yl{Gaussian} splatting and define the texture parameters we use in rendering (Sec.~\ref{subsec:GSR}).  
Since directly optimizing everything from scratch can lead to suboptimal results, especially on the geometry side, we introduce a normal field distillation module (Sec.~\ref{subsec:NFD}) that utilizes an SDF network to guide \yl{Gaussians} to a better surface representation. 
With geometry set up properly, we bind texture parameters onto the \yl{Gaussian} splatting representation and perform \textit{deferred shading} to obtain plausible decomposition (Sec.~\ref{subsec:DS_GS}). 
Finally, we demonstrate how this decomposition enables scene editing applications like relighting and texture editing (Sec.~\ref{subsec:Edit}). 

\begin{figure*}[t!]
	\includegraphics[width=0.99\linewidth]{./img/pipeline_toaster.pdf}
	\caption{
		Overview of \name. 
        Given multi-view images, \name builds a decoupled \yl{Gaussian} splatting representation where auxiliary texture attributes are defined on each \yl{Gaussian}. 
        To enable successful geometry and appearance separation, we attach an extra normal direction to each \yl{Gaussian} and optimize it by distilling the normal field from a jointly trained signed distance function. 
        For texture and lighting decomposition, \name rasterizes geometry and texture attributes into buffer maps and computes shading at \yl{the} pixel level under the illumination of a learnable environment map with \textit{deferred shading}. 
	}
	\label{fig:pipeline}
\end{figure*}



\subsection{Gaussian Splatting and Texture Definition}
\label{subsec:GSR}
Gaussian splatting models a 3D scene by a set of \yl{Gaussians} with which attributes like position $P$, 
covariance $\Sigma$, 
scale $S$, opacity $\alpha$ and Spherical Harmonic coefficients (SH) representing view-dependent appearance are associated. 
It uses a point-based $\alpha$-blending rendering formulation similar to NeRF~\cite{NeRF}. 
\begin{equation}
\label{eqn:GSR}
C_* = \sum_{i=1}^N *_i \alpha_i T_i 
\end{equation}
where $N$ is the number of sample points on a ray. 
$T_i = \prod_{j=1}^{i-1} (1-\alpha_j)$ is accumulated transmittance and $*$ can be any attribute of the sample point to be blended such \yl{as} depth, color and normal direction. 
Gaussian splatting avoids densely sampling on the ray and efficiently projects 3D \yl{Gaussians} onto a 2D image plane and approximates 3D scenes by blending view-dependent color represented by SH. 
Instead of modeling the appearance as \yl{a} whole with SH, \yl{our} \name replaces SH with optimizable texture parameters including diffuse albedo $k_d$, roughness $r$, and specular albedo $k_s$ as shown in Fig.~\ref{fig:pipeline} to enable texture and lighting decoupling. 

\subsection{Normal Field Distillation}
\label{subsec:NFD}

Normal estimation is important for extracting correct texture and lighting information. 
However, it is challenging to obtain plausible surface \yl{normals} in \yl{Gaussian} splatting as mentioned in Sec.~\ref{sec:intro} since it is a discrete representation and only promotes rendering quality without constraining the geometry. 
To resolve this issue, we propose a normal field distillation module that integrates current multi-view geometry reconstruction techniques with \yl{Gaussian} splatting. 
Specifically, we jointly train a NeRF-like network equipped with multi-level hash encoding~\cite{InstantNGP} namely \neusbranch 
and a \yl{Gaussian} splatting representation. 
For the \neusbranch, we utilize the signed distance field (SDF) as the geometry representation which can be parameterized as a neural network $f_s(x)$ where $x$ is a sample point in 3D space. 
For the color branch, instead of modeling all the appearances with a single view-dependent network like NeRF~\cite{NeRF}, we separate the appearance into a view-independent branch and a view-dependent branch to avoid geometric artifacts like concave surfaces on reflective scenes~\cite{Ref-NeuS, Factored-NeuS}. 
The view-independent branch predicts the diffuse color $c_d$ and the specular tint $p$ for a sample point $x$. 
The view-dependent branch takes the view direction as input and outputs the specular lighting $c_l$. 
The final color can be composed by $c' = c_d + p c_l$. 
For better understanding, all terms with superscript $'$ refer to the NeRF network while terms without refer to the Gaussian splatting representation. 
To render a pixel $C_c'(v)$, we integrate colors of sample points on a ray $v$ with Eqn.~(\ref{eqn:GSR}).
Adopting the occlusion-aware and unbiased volume rendering technique from NeuS~\cite{NeuS}, the SDF value can be transformed to volume density by $\sigma(t) = \max\left(-\frac{\frac{d\Phi_s}{dt}(f(x(t))}{\Phi_s(f(x(t))}, 0\right)$, where $\mathit{\Phi_s(x)=(1+e^{-sx})^{-1}}$ and $\mathit{s}$ is a trainable deviation parameter.
The opaque value for point $x_i$ is derived from the density function $\sigma(t)$ by $\alpha_i = 1 - exp\left(-\int_{t_i}^{t_{i+1}} \sigma(t) dt\right)$. 
The loss function for \neusbranch is as follows:
\begin{equation}
    \label{eqn:loss_refneus}
    L_{nerf} = \sum_{v {\in V}} \left\|C_c'(v) - C^{t}(v) \right\| + \lambda \sum_{v {\in V}} \sum_{i{=1}}^{M} \left\| \|\nabla_{x_{v, i}} \| - 1 \right\|_2^2,
\end{equation}
where $V$ is the number of rays in a batch. 
$M$ is the number of sample points on a single ray and $C^{t}(v)$ is the corresponding ground truth color. 
The second term is the Eikonal loss, where $\|\nabla f_s({x_{v, i})} \|$ is the spatial norm of the SDF network $f_s(x)$'s gradient at \yl{the} $i$th sample point $x_{v, i}$ on the ray $v$.
In order to distill the normal information from \neusbranch to \yl{Gaussian} splatting, we define an extra normal direction $n$ on each \yl{Gaussian} as shown in Fig.~\ref{fig:pipeline} and render the normal $C_n$ of Gaussian splatting for a pixel using Eqn.~(\ref{eqn:GSR}). 
Then we distill the normal field \yl{Gaussian} splatting by minimizing the difference between $C_n$ and the corresponding rendered normal in \neusbranch:
\begin{equation}
\label{eqn:L_nd}
\begin{aligned}
     L_{nd} = \sum_{v \in V} \left\|1 - C_n(v) \cdot C_n'(v)\right\| = \sum_{v \in V} \left\|1 - C_n(v) \cdot C_{\nabla x_{v,i}} \right\|
\end{aligned}
\end{equation}
where $\cdot$ is the dot product operator. 

\subsection{Deferred Shading in Gaussian Splatting}
\label{subsec:DS_GS}
For shading calculation, \yl{we} follow the rendering equation~\cite{RenderingEquation}:
\begin{equation}
\label{eqn:rendering_eqn}
    L_o(\omega_{o}) = \int_{\Omega} L_i(\omega_{i}) f(\omega_{i}, \omega_{o}) (\omega_{i} \cdot n) d \omega_i
\end{equation}
where $L_o(\omega_{o})$ and $L_i(\omega_{i})$ represent outgoing lighting in direction $\omega_{o}$ and incident lighting from direction $\omega_{i}$ respectively. 
$f(\omega_{i}, \omega_{o})$ is the BRDF 
at a point which can be parameterized by the Disney shading model~\cite{DisneyBRDF}:
\begin{equation}
    f(\omega_{i}, \omega_{o}) = \frac{k_d}{\pi} + \frac{DFG}{4(\omega_{i} \cdot n)(\omega_{o} \cdot n)}
\end{equation}
where $D, F, G$ correspond to normal distribution function, \yl{Fresnel} term, and geometry term. 
Their detailed calculations can be found in~\cite{microfacet}. 

We model the scene's environment lighting with a $6 \times 512\times 512$ High Dynamic Range (HDR) cube map and utilize the SplitSum~\cite{SplitSum} approximation which separates the integrals of lighting and BRDF to enable efficient shading calculation. 
T   he diffuse color $c_{diff}$:
\begin{equation}
    c_{diff} = \frac{k_d}{\pi} \int_{\Omega} L_i(\omega_{i}) (\omega_{i} \cdot n) d \omega_i
\end{equation}
which is only dependent on the normal direction $n$ and can be pre-computed and stored in a 2D texture. 

For the specular color $c_{spec}$:
\begin{equation}
\begin{aligned}
    &c_{spec} \approx Int_{light} \cdot Int_{BRDF} \\
            &= \int_{\Omega}  L_i(\omega_{i}) D(\omega_i, \omega_o) (\omega_{i} \cdot n) d \omega_i \cdot
    \int_{\Omega} f(\omega_{i}, \omega_{o}) (\omega_{i} \cdot n) d \omega_i
\end{aligned}
\end{equation}
where $Int_{light}$ represents the integral of incident light with the normal distribution function $D$. 
Given a fixed environment map, $Int_{light}$ is only related to the roughness value $r$ and therefore can be pre-computed and stored in a mipmap where each mip level corresponds to a fixed roughness value. 
$Int_{BRDF}$ is the integral of specular BRDF under a uniform white environment lighting. 
It is determined by the roughness value in the BRDF and the dot product between incident light direction and normal direction $\omega_{i} \cdot n$. 
Again it can be pre-computed and stored in a 2D texture and efficiently looked up in rendering. 
The final shaded color is $c_g = c_{diff} + c_{spec}$. 

A straightforward way to obtain a decoupled \yl{Gaussian} splatting representation is performing \textit{forward shading} for each \yl{Gaussian} to get its color $c_i$ and rasterizing \yl{Gaussians} with Eqn.~(\ref{eqn:GSR}) similar to previous works~\cite{invrender, NeRFactor, GIR, GaussianShader}. 
Though obtaining plausible decoupled results, we observe notable artifacts when \yl{rendering} the decoupled \yl{Gaussians} under a novel lighting condition (see ``F.S." columns in Fig.~\ref{fig:ablate_DS} and ``RGS" and ``GShader" columns in Fig.~\ref{fig:relight}). 
This is because even \yl{though} Eqn.~(\ref{eqn:L_nd}) constrains the \yl{rasterized} normal to be close to that of \neusbranch at pixel level, the \yl{rasterized} normal is still blended from multiple \yl{Gaussians}' normals. 
The geometry attributes like position, covariance and opacity may overfit the given lighting condition to produce proper blended normal and shaded color and it does not guarantee that each \yl{Gaussian}'s individual normal or shaded color is correct. 
Thus when changing the lighting condition, the trained geometry attributes may not work for the shaded colors under a new lighting condition and blend them in \yl{a} wrong way.  

To resolve this issue, we introduce the \textit{deferred shading} technique into the rendering of \yl{Gaussians}. 
That is to say, we first rasterize components like position $P$, normal $n$, diffuse albedo $k_d$, roughness $r$, specular albedo $k_s$, and opacity $\alpha$ into 2D pixels $C_{P}, C_{n}, C_{k_d}, C_{r}, C_{k_s}$, and $C_{\alpha}$ with Eqn.~(\ref{eqn:GSR}). 
By aggregating all pixels, we get corresponding 2D maps $I_{P}, I_{n}, I_{k_d}$, $I_{r}, I_{k_s}$, and $I_{\alpha}$. 
Then we calculate the shaded color $C$ for a pixel with Eqn.~(\ref{eqn:rendering_eqn}). 
Note that since computing shading for every single pixel in the image is time-consuming, we only perform shading on pixels with $C_{\alpha}$ bigger than 0.5. 
Overall, we optimize \yl{the} following loss:
\begin{equation}
\label{eqn:total_loss}
\begin{aligned}
    L = &L_{nerf} + \lambda_{nd} L_{nd} + \lambda_{L1}L_{L1} + \lambda_{ssim}L_{ssim} +\\
    & \lambda_{mask}L_{mask} + \lambda_{TV} L_{TV}
\end{aligned}
\end{equation}
where $L_{L1}$ and $L_{ssim}$ are the L1 loss and D-SSIM loss between the rasterized image and the ground truth image same as those in~\cite{GS}. 
$L_{mask} = \|I_{\alpha} - I_m\|_1$ is the mask loss that encourages \yl{Gaussians} to lie in the foreground and $I_m$ is the ground truth mask image. 
$L_{TV}$ is the total variation loss applied on images $I_{k_d}, I_{k_s}, I_{k_s}$ to encourage smooth texture estimation. 



\subsection{Editing \yl{Decoupled} Gaussians}
\label{subsec:Edit}
After optimizing the attributes on \yl{Gaussians}, we obtain a decoupled \yl{Gaussian} splatting representation and can render the original scene with edited texture or geometry under a novel illumination. 
\subsection{Relighting}
As we extract the environment map of the original scene, we can replace it by a novel environment map \yl{with} little effort and render the scene with novel illumination using \textit{deferred shading}. 

\subsection{Geometry Editing}
For geometry editing, instead of directly operating on \yl{Gaussians}, we first extract a mesh with \neusbranch and use it as a geometry proxy to guide the deformation of \yl{Gaussians}. 
Specifically, we first deform the mesh using the \yl{As-Rigid-As-Possible (ARAP)} deformation algorithm~\cite{ARAP} which solves \yl{for} rotations and translations of mesh vertices.  
Then we find the closest neighbor 
for each \yl{Gaussian} on the original mesh \yl{surface}, whose rotation and translation can be obtained via barycentric mapping \yl{(based on the mesh triangle that contains the neighboring point)}. 
Finally, we deform each \yl{Gaussian} by applying the rotation and the translation of its corresponding closest neighbor to \yl{Gaussian}'s position, covariance matrix, and normal direction.
By rasterizing deformed \yl{Gaussians}, we can render the deformed scene. 

\subsection{Texture Editing}
\label{sec:texture_edit}
For texture editing, we follow the editing paradigm in traditional 3D modeling software~\cite{blender} to paint the scene from a given viewpoint. 
Our texture editing supports changes on all texture components. 
We first determine which Gaussians need to be adjusted by considering both whether they lie in the editing mask and the difference between the \yl{Gaussian} depth and the corresponding depth in \yl{the} depth map rendered with Eqn.~(\ref{eqn:GSR}). 
Then we can optimize these \yl{Gaussians'} texture attributes $*$ by:
\begin{equation}
\label{eqn:edit_naive}
    {\arg}\min_{*} \|I_{*}^i - I_e^i\|,\space\space\space * \in \{k_d, r, k_s\}
\end{equation}
where $I_{*}^i$ denotes a rendered texture map, \yl{e.g.} diffuse albedo map, from the input editing viewpoint $i$. 
$I_e^i$ is the edited texture map from viewpoint $i$ by putting edited \yl{content} onto $I_{*}^i$. 

However, this naive solution only considers the input editing viewpoint and may produce blending artifacts when we view the optimized \yl{Gaussians} from other viewpoints. 
Therefore, we propose to augment the optimization with random view inputs. 
That is to say, we generate edited rendered results from multiple random viewpoints by first \yl{unprojecting the} edited \yl{content} onto the mesh extracted by \neusbranch to form a textured mesh and \yl{rendering} the textured mesh from random viewpoints. 
So the final editing loss is: 
\begin{equation}
\label{eqn:edit_improved}
    {\arg}\min_{*} \sum_{j \in L}\|I_{*}^j - I_e^j\|,\space\space\space * \in \{k_d, r, k_s\}
\end{equation}
where set $L$ includes both the input viewpoint and randomly sampled viewpoints.

